\documentclass[11pt]{article}

\usepackage[margin=1in]{geometry}
\usepackage{graphicx}
\usepackage{array}
\usepackage{hyperref}
\usepackage{enumitem}
\usepackage{booktabs}
\usepackage{caption}
\usepackage{amsmath}
\usepackage{titlesec}

\titleformat{\section}{\large\bfseries}{\thesection.}{1em}{}

\begin{document}

\begin{center}
    \Large \textbf{Sri Sivasubramaniya Nadar College of Engineering, Chennai} \\
    \large (An Autonomous Institution affiliated to Anna University) \\
    \vspace{0.3cm}
\end{center}

\begin{table}[h!]
\renewcommand{\arraystretch}{1.5}
\centering
\resizebox{\textwidth}{!}{%
\begin{tabular}{|l|cll|}
\hline
\textbf{Degree \& Branch} & \multicolumn{1}{c|}{B.E Computer Science \& Engineering} & \textbf{Semester} & VI \\ \hline
\textbf{Subject Code \& Name} & \multicolumn{3}{c|}{UCS2612 -- Machine Learning Laboratory} \\ \hline
\textbf{Academic Year} & \multicolumn{1}{c|}{2025--2026 (Even)} & \textbf{Batch} & 2023--2027 \\ \hline
\end{tabular}
}
\end{table}

\begin{center}
    \textbf{\Large Experiment 7}
\end{center}

\begin{center}
    \textbf{Bagging, Boosting, and Stacked Ensemble Models}
\end{center}

%--------------------------------------------------

\section*{Objective}
\begin{itemize}
    \item To understand ensemble learning strategies: Bagging, Boosting, and Stacking.
    \item To implement Bagging and Boosting classifiers.
    \item To build a Stacked Ensemble model using multiple base learners.
    \item To compare ensemble models in terms of accuracy, stability, and generalization.
    \item To analyze the effect of ensemble methods on bias and variance.
\end{itemize}

%--------------------------------------------------

\section*{Dataset}
\textbf{Wisconsin Diagnostic Breast Cancer Dataset}

\begin{itemize}
    \item Total samples: 569
    \item Features: 30 numerical attributes
    \item Target classes: Malignant (M) and Benign (B)
\end{itemize}

\textbf{Dataset Link:}  
\href{https://archive.ics.uci.edu/dataset/17/breast+cancer+wisconsin+diagnostic}{https://archive.ics.uci.edu}

%--------------------------------------------------

\section*{Theory}

\subsection*{Bagging (Bootstrap Aggregation)}
Bagging is an ensemble technique that trains multiple models on different bootstrap samples of the training data and aggregates their predictions.

\textbf{Key Characteristics:}
\begin{itemize}
    \item Reduces variance
    \item Suitable for high-variance models
    \item Independent model training
\end{itemize}

\subsection*{Boosting}
Boosting is an ensemble method where models are trained sequentially, and each new model focuses on correcting the errors made by previous models.

\textbf{Key Characteristics:}
\begin{itemize}
    \item Reduces bias
    \item Sequential learning
    \item Emphasizes difficult samples
\end{itemize}

Common boosting algorithms include AdaBoost and Gradient Boosting.

\subsection*{Stacked Ensemble (Stacking)}
Stacking combines multiple heterogeneous base models and trains a meta-learner to optimally combine their predictions.

\textbf{Key Characteristics:}
\begin{itemize}
    \item Uses diverse base learners
    \item Meta-model learns optimal combination
    \item Often achieves superior performance
\end{itemize}

%--------------------------------------------------

\section*{Steps for Implementation}

\begin{enumerate}[label=\arabic*.]
    \item Load and preprocess the dataset.
    \item Perform Exploratory Data Analysis (EDA).
    \item Split the dataset into training and testing sets (80--20).
    \item Implement a Bagging classifier using a base estimator.
    \item Implement Boosting classifiers (AdaBoost and/or Gradient Boosting).
    \item Construct a Stacked Ensemble using multiple base models.
    \item Define hyperparameter search spaces for each ensemble model.
    \item Select hyperparameters using 5-fold cross-validation.
    \item Evaluate all models using standard performance metrics.
    \item Compare results and analyze bias--variance behavior.
\end{enumerate}

%--------------------------------------------------

\section*{Ensemble Models Used}

\subsection*{Bagging}
\begin{itemize}
    \item Base Estimator: Decision Tree
    \item Number of estimators
    \item Sampling strategy
\end{itemize}

\subsection*{Boosting}
\begin{itemize}
    \item AdaBoost
    \item Gradient Boosting
\end{itemize}

\subsection*{Stacked Ensemble}
\begin{itemize}
    \item Base Models: SVM, Na\"ive Bayes, Decision Tree
    \item Meta Learner: Logistic Regression
\end{itemize}

%--------------------------------------------------

\section*{Hyperparameters to be Explored}

\subsection*{Bagging}
\begin{itemize}
    \item \texttt{n\_estimators}
    \item \texttt{max\_samples}
    \item \texttt{max\_features}
\end{itemize}

\subsection*{Boosting}
\begin{itemize}
    \item \texttt{n\_estimators}
    \item \texttt{learning\_rate}
    \item \texttt{max\_depth}
\end{itemize}

\subsection*{Stacked Ensemble}
\begin{itemize}
    \item Choice of base models
    \item Final estimator
\end{itemize}

%--------------------------------------------------

\section*{Hyperparameter Evaluation Results}

\subsection*{Bagging Results}
\begin{table}[h!]
\centering
\caption{Bagging Hyperparameter Evaluation}
\begin{tabular}{cccc}
\toprule
n\_estimators & max\_samples & Avg CV Accuracy (\%) & Avg CV F1 Score \\
\midrule
 &  &  &  \\
 &  &  &  \\
 &  &  &  \\
\bottomrule
\end{tabular}
\end{table}

\subsection*{Boosting Results}
\begin{table}[h!]
\centering
\caption{Boosting Hyperparameter Evaluation}
\begin{tabular}{cccc}
\toprule
n\_estimators & learning\_rate & Avg CV Accuracy (\%) & Avg CV F1 Score \\
\midrule
 &  &  &  \\
 &  &  &  \\
 &  &  &  \\
\bottomrule
\end{tabular}
\end{table}

\subsection*{Stacked Ensemble Results}
\begin{table}[h!]
\centering
\caption{Stacked Ensemble Evaluation}
\begin{tabular}{lccc}
\toprule
Base Models & Meta Learner & Avg CV Accuracy (\%) & Avg CV F1 Score \\
\midrule
 &  &  &  \\
 &  &  &  \\
\bottomrule
\end{tabular}
\end{table}

%--------------------------------------------------

\section*{Performance Comparison}

\begin{table}[h!]
\centering
\caption{Performance Comparison of Ensemble Models}
\begin{tabular}{lcccc}
\toprule
Model & Accuracy (\%) & Precision & Recall & F1 Score \\
\midrule
Bagging &  &  &  &  \\
Boosting &  &  &  &  \\
Stacked Ensemble &  &  &  &  \\
\bottomrule
\end{tabular}
\end{table}

%--------------------------------------------------

\section*{Evaluation Metrics}
\begin{itemize}
    \item Accuracy
    \item Precision
    \item Recall
    \item F1-score
    \item Confusion Matrix
    \item ROC Curve and AUC
\end{itemize}

%--------------------------------------------------

\section*{Observation Questions}
\begin{itemize}
    \item How does Bagging reduce variance?
    \item How does Boosting address model bias?
    \item Why does stacking benefit from heterogeneous models?
    \item Which ensemble method performed best and why?
\end{itemize}

%--------------------------------------------------

\section*{Conclusion}
Bagging, Boosting, and Stacked Ensemble models were implemented and evaluated.  
The experiment demonstrates how different ensemble strategies improve predictive performance and generalization compared to individual models.

%--------------------------------------------------

\section*{References}
\begin{itemize}
    \item \href{https://scikit-learn.org/stable/modules/ensemble.html}{Scikit-learn: Ensemble Methods}
    \item \href{https://scikit-learn.org/stable/modules/generated/sklearn.ensemble.BaggingClassifier.html}{Bagging Classifier}
    \item \href{https://archive.ics.uci.edu/dataset/17/breast+cancer+wisconsin+diagnostic}{UCI Dataset}
\end{itemize}

\end{document}
